\documentclass[
    headings=optiontohead,              % allows double headers
    12pt,                               % fontsize 
    DIV=13,                             % koma script diveider amount. tells koma how much of the site can be written to
    twoside=false,                      % if set to true, automatically formats as book style with different left and right pages
    open=right,                         % starting page on twosided texts 
    BCOR=00mm,                          % correction that accounts for the center of the pages being glued in
    toc=bibliographynumbered            % bibliography gets a number and is listed in the table of contents
]{scrreport}

\usepackage[english]{babel}                     % font that supports English
\usepackage{upgreek}                            % non-cursive Greek letters
\usepackage[stretch=10,shrink=10,protrusion=true,expansion=true,final]{microtype} % prettier block format
\usepackage{hyperref}                           % links for everything
\usepackage{color}                              % allows for setting in different colors
\usepackage[autooneside=false,automark]{scrlayer-scrpage} % page-style with "Kolumnentitel" (title of current chapter is displayed at the top)
\usepackage[sb]{libertinus}                     % use the font libertinus (needs to be installed from the web)
\usepackage[slantedGreek]{libertinust1math}     % math mode improvement for libertinus
\usepackage{siunitx}                            % physical units setting
\usepackage{icomma}                             % commas in lists get extra space if needed                        
\usepackage{amsfonts,amstext,amsmath,amsthm} % better math mode (\mathrm and \text) and symbols
\usepackage{xspace}                             % works to improve own commands and provides "\xspace"-command, that puts a space if needed
\usepackage{ifthen}                             % more control over non-obligatory parameters
\usepackage{titling}                            % get title values as macros
\usepackage[onehalfspacing]{setspace}           % control the spacing between lines and in enumeration lists
\usepackage[backend=biber, style=phys, biblabel=brackets]{biblatex} % citations with "modern" backend and an physics-accepted citation style
\usepackage{graphicx}                           % work with graphics 
\usepackage{ragged2e}                           % ragged-commands (when no block format is wanted)
\usepackage{pdfpages}                           % allows including of pdfs into this pdf
\usepackage{booktabs}                           % better table formatting
\usepackage{multicol}                           % allows for the definition of multi-columns in tables
\usepackage{multirow}                           % allows for the definition of multirow-tables instead of just multicolumn
\usepackage[section]{placeins}                  % provides the command "\FloatBarrier" to control the end of floatable regions for figures/tables
\usepackage{float}                              % provides the "H" option for forcing placement of a figure
\usepackage{floatpag}                           % make it possible for float-pages to not have a page number
\usepackage{url}                                % sometimes needed by biblatex, technically no longer needed
\usepackage{minted}                             % nice code highlighting (needs Python Package to compile!!)
\usepackage{accents}                            % better control over accents
\usepackage{mathtools}                          % more math control possibilities
\usepackage[autostyle=true]{csquotes}           % context-sensitive-quotes -> quotation marks that are set correctly for the context
\usepackage{physics}                            % bra-ket and more
\usepackage{nicematrix}                         % label row/cols on matrix
\usepackage{caption}                            % caption of different environments

\title{Measurement of Droplets in vapourised fluids using Machine Learning Techniques}
\author{Jan Claar}
\date{\today}

\graphicspath{{./images/}} % custom paths for folders in that graphics can befound

\sisetup{ % setup for siunitx  
    detect-all, 
    locale=US, % language setup for siunitx 
    range-phrase={ \text{to} }, % word that is put into an si range
    range-units = single, % better display of error ranges
    per-mode=symbol-or-fraction, % more dynamic frac usage in inline/displaymathmode 
    separate-uncertainty, % for better +- , \pm when including an error range 
}

\hypersetup { 
    colorlinks=true, linkcolor=dblue, % dark blue linkcolor
    urlcolor=dblue, % dark blue linkcolor citecolor=dblue, % dark blue linkcolor
    pdfauthor = {Jan Claar}, % write details into the expanded file properties
    pdftitle = {Investigation of transformer architectures for geometrical graphstructures}, 
    pdfkeywords = {neural networks, ai, physics, machine learning}, pdfsubject = {Bachelor Thesis} 
}

\captionsetup[sub]{font=small}

\AtBeginDocument{ 
    \let\mathbb\relax
    \DeclareMathAlphabet\PazoBB{U}{fplmbb}{m}{n} \newcommand{\mathbb}{\PazoBB} 
}
%more options to the \mathbb command

\setminted[]{ 
    xleftmargin=0cm, 
    xrightmargin=0cm, 
    frame=single, 
    framesep=.25cm,
    linenos, 
    tabsize=2, 
    breaklines, 
    breakafter=.], 
    breakaftersymbolpre= , 
}
%configure the minted code-highlighting style

\addbibresource{literature.bib} %initialize bibtex with correct file

\NiceMatrixOptions{ 
    code-for-first-row = \color{dblue} , 
    code-for-last-row =\color{dblue} , 
    code-for-first-col = \color{dblue} , 
    code-for-last-col = \color{dblue} 
}                                  % another file that holds the package/document configuration
\linespread{1.1} % line-spacing can be controlled here

\clubpenalty10000 % Schusterjunge, orphan 
\widowpenalty10000 % Hurenkind, Witwe
\displaywidowpenalty=10000 % Make document obey stricter rules considering "Schusterjungen" and "Hurenkinder" \setcounter{biburlnumpenalty}{100}
\setcounter{biburlucpenalty}{100} 
\setcounter{biburllcpenalty}{100}
\renewcommand{\topfraction}{0.8} % allows for more chilled "text to image ratio"
\renewcommand{\bottomfraction}{0.8} 
\renewcommand{\textfraction}{0.1}
\renewcommand{\floatpagefraction}{0.8}

% \renewcaptionname{ngerman}{\figurename}{Abb.} %"Figure" becomes "Fig." inE nglish 
\setcapindent{0cm} %useful if image captions have multiple lines. Removers indentation below "Fig." 
\setlength{\parindent}{0cm} %removes indentation at start of new paragraphs 
\setlength{\abovecaptionskip}{0.0cm}

%bibliography slots are redefined/modified here
\DeclareFieldFormat{journaltitle}{\textsl{#1}\isdot}
\DeclareFieldFormat{titlecase}{{#1}}

%COLORS 
\definecolor{dblue}{rgb}{0,0,0.5} 
\definecolor{dred}{rgb}{0.5,0,0}
\definecolor{dgrey}{rgb}{0.5,0.5,0.5} 
\definecolor{lgrey}{rgb}{0.8,0.8,0.8}
\definecolor{textred}{RGB}{170,0,0} 
\definecolor{textyellow}{RGB}{170,170,0}
\definecolor{textgreen}{RGB}{0,170,0} 
\definecolor{textblue}{RGB}{0,0,170}

%overwrite the coma-script definitions
\addtokomafont{pagehead}{\normalfont\color{dgrey}} %overwrite the coma-scriptdefinitions 
\addtokomafont{sectioning}{\rmfamily\color{dblue}\boldmath}
%rmfamily puts headings in "normal" "serif-font" instead of "sans-serif" boldmath ensures a bold math font in subscripts
\addtokomafont{captionlabel}{\bfseries\footnotesize} %better Fig. format
\addtokomafont{caption}{\footnotesize}

%! hyphenation commonly used words can be spelled here to provide latex with the correct places to make line breaks 
\hyphenation{Li-pid-mono-lage}

% headline spacing 
\RedeclareSectionCommand[beforeskip=0cm, afterskip=1cm]{chapter}                                  % another file that holds format information
%! Ref-Commands 
\newcommand*{\fullref}[1]{\hyperref[{#1}]{\textit{\autoref*{#1}\nameref*{#1}}}} 
\newcommand*{\fullpage}[1]{\hyperref[{#1}]{Seite\pageref*{#1}}} 
\newcommand*{\fullpages}[1]{\hyperref[{#1}]{Seiten\pageref*{#1}ff}}

%! Math operators and other small conveniences
\newcommand\thickbar[1]{\accentset{\rule{.6em}{.8pt}}{#1}}
\DeclareMathOperator{\ggt}{ggT} \DeclareMathOperator{\kgv}{kgV}
\DeclarePairedDelimiter\ceil{\lceil}{\rceil}
\DeclarePairedDelimiter\floor{\lfloor}{\rfloor} 
\renewcommand*{\arraystretch}{0.8}
\DeclareSIUnit\pixel{px}

\newcommand{\filepath}[2]{\colorbox{lgrey}{#1: \texttt{#2}}\xspace}

%! replace German quotation marks with English ones (I always write the left and right quotes explicitly with the German command variants as they are nice to debug and already in my macros. Then in the end, I can swap them with the desired replacement quotes of my choice) 
\renewcommand*{\glqq}{``}
\renewcommand*{\grqq}{''}                                % another file that holds predefined commands

\begin{document}
    \thispagestyle{empty} 
    \newcommand{\mail}{jan.claar@student.uni-augsburg.de}

\begin{titlepage} 
    \makebox[\textwidth][c]{\includegraphics[width=0.5\textwidth]{images/logo_uni_augsburg.jpg}} \color{dblue}
    \begin{center} \vspace*{2cm} \Huge \textbf{\thetitle}
        
        \vspace*{1.5cm} 
        
        \color{black} \textbf{Bachelor Thesis}

        \vspace*{1cm} 

        \normalsize submitted by\\ \LARGE \theauthor\\\vspace*{0.3cm} \normalsize on \thedate

        \vspace{1.8cm} 

        \color{black} \emph{Augsburg University}\\ \emph{Faculty ofApplied Computer Science}\\ \emph{Institute of Computer Science}\\ \emph{Chair for Machine Learning \& Computer Vision}
        \vfill

        \begin{tabular}{rl}
            Supervisor: & Daniel Kienzle\\ 
            1$^\text{st}$ Corrector: &Prof. Dr. Rainer Lienhart\\
            2$^\text{nd}$ Corrector: &Prof. Dr. \\% TODO 
        \end{tabular} 
    \end{center}

\end{titlepage} 
    \tableofcontents
    \thispagestyle{empty} 
    \newpage 
    \setcounter{page}{1} 
    \pagestyle{scrheadings}

    \chapter{Introduction}
    \label{sec:introduction}
    The advancement of all major natural sciences has often gone hand in hand with advances in the respective measurement techniques used to formulate and validate hypotheses through experiments using the scientific method. 
Being able to measure key metrics in the system of interest is vital not only in scientific research, but also the developement of new technologies and products. 

The field of \emph{machine learning} and \emph{artificial intelligence} has seen an explosive growth over the past two decades and has become a major research focus for many in the discipline of data and computer science. 
With storage capacity and processing speed of computer hardware, such as specialized \textbf{G}raphics and \textbf{T}ensor \textbf{P}rocessing \textbf{U}nits (GPUs/ TPUs) experiencing similar advancements, the times are long gone when using AI to solve a problem was almost always unreasonable because of its large demand for data and computational resources. 

It is no wonder, then, that we ask how we can apply the power of machine learning to the field of measurement science. In fact, this question was asked even before the 2000s \cite{alippiArtificialIntelligenceInstruments1998}, at time when access to AI was much more restricted. Since then AI has been and will continue to be successfully used to improve measurement techniques by helping develop better sensors \cite{ballardMachineLearningComputationenabled2021}, processing raw sensor data or deriving meaning from already processed data. 

One way in which machine learning (in this case neural networks) can be used to improve the measuring process is by automating tasks that are difficult to solve through classical algorithms, but which are easy for a human to do. This is no surprise, since the idea of artificial neural networks in the first place is inspired by their biological counterpart. These problems often are very intuitively solved by the brain for a single instance, but require a large amount of effort when scaled. 
It is this aspect of machine learning that we hope to employ in our research. 

\textbf{S}urface \textbf{A}coustic \textbf{W}aves (SAW) can be produced by using electrodes to induce surface vibrations in piezoelectric substrates. These SAW chips have several useful applications, one of them being using them to vapourise fluids into very small droplets, which could be beneficial in fields such as e.g. medicine, for producing fine vapour of solutions of certain drugs to ensure optimal absorption in the body.
While researching and developing this technology it is important to obtain insight about the characteristics of the produced vapour, primarily droplet size and distribution. While there are already techniques to measure the size of droplets in a vapour, they all come with certain caveats. A simple solution is then, to take high speed image data of the vapour and measure the droplet size directly on the images, which is easy for one image, but constitutes an ardous task for each measurement run, if one desires to obtain any statistically robust data.

The goal of this thesis is, now, to explore the use of neural networks to automatically identify droplets in image data taken of SAW vapourised fluids and measure them with sufficient precision, as well as to compare this approach to using classical algorithms for image segmentation. Finally, a simple application will be developed which should find use in current research.

    \newpage 
%\nocite{*} 
\printbibliography[title={Bibliography}]

\end{document}